% !TEX root = 00_dmlr_compel.tex
% New section for the DMLR version. Reviewers of an earlier draft judged the method
% "conceptually simple" and asked for deeper theoretical grounding; this section gives
% the Shannon / Kolmogorov / MDL reading of the Goldilocks band.
\section{Theoretical Grounding: Compression, Entropy, and Learnable Structure}\label{sec:theory}

A natural objection to \textsc{Compel} is that a single scalar seems too simple a statistic to carry information about data quality.
Information theory suggests otherwise: compression ratio estimates a document's entropy rate, and entropy rate separates the three regimes visible in Figure~\ref{fig:examples}.
This section makes that argument explicit.

\paragraph{Compression ratio estimates entropy rate.}
Shannon's source coding theorem states that no lossless code can use fewer expected bits per symbol than the entropy rate of the source, and that good codes approach this bound \citep{shannon1948mathematical}.
A practical compressor therefore acts as an entropy estimator: the length of $c(x)$ upper-bounds the information content of $x$ under the compressor's model class.
Shannon himself estimated the entropy of printed English this way, by measuring how far text could be predicted and hence compressed \citep{shannon1951prediction}.
In algorithmic terms, the ideal compressed length is the Kolmogorov complexity $K(x)$---the length of the shortest program that outputs $x$ \citep{kolmogorov1965three}.
$K$ is uncomputable, and real compressors are its standard computable surrogates \citep{li2008introduction}.
The same substitution underlies the Normalized Compression Distance \citep{cilibrasi2005clustering} and ZIP-FIT \citep{obbad2024zipfit}; $\mathrm{CR}(x)$ in Equation~\eqref{eq:cr} is simply the per-byte, normalized form of this surrogate.

\paragraph{Language modeling is compression.}
A language model with next-token distribution $p$ defines a lossless code that assigns $x$ a code length of $-\log_2 p(x)$ bits, so minimizing pretraining loss is exactly minimizing description length \citep{deletang2023language}.
Empirically, compression capacity tracks model capability \citep{huang2024compression}, and gzip compressibility predicts how scaling laws shift with data composition \citep{pandey2024gzip}.
Pretraining data value is therefore naturally measured in bits: a document is useful to the extent that it contains structure the model does not yet encode but can learn to encode.

\paragraph{Why both extremes are poor training data.}
The minimum description length (MDL) principle formalizes learning as two-part compression: a model part that captures regularity, and a residual part for what the model cannot explain \citep{rissanen1978modeling, grunwald2007minimum}.
The two tails of the compression-ratio distribution fail this decomposition in opposite ways.
When $\mathrm{CR}(x)$ is far below the band, the document's entropy rate is small: after a short prefix, additional tokens contribute almost no new bits (templated listings, keyword stuffing, boilerplate).
Gradient updates on such tokens repeat what the model has already absorbed, so we conjecture that they buy little generalization per FLOP.
When $\mathrm{CR}(x)$ is far above the band, the document is nearly incompressible for the LZ family, i.e., statistically close to random relative to that model class (hexadecimal tables, encoding artifacts, shuffled mixed-language spam).
A string that is algorithmically random admits no shorter description at all \citep{li2008introduction}; its loss is an irreducible noise floor, and capacity spent fitting it is capacity spent memorizing noise.

\paragraph{The interior band is where learnable structure lives.}
Between the extremes lies text with regularities that a weak universal compressor captures only partially: LZ4 exploits local repetition but not syntax, discourse, or semantics.
Precisely this residual structure is what a higher-capacity learner can still compress---that is, learn.
Well-edited human prose empirically concentrates in this regime (Figures~\ref{fig:cr-distributions} and~\ref{fig:fwedu-kde}), which explains why an interior compression-ratio band recovers much of the behavior of a trained quality classifier at a vanishing fraction of its cost.

\paragraph{Caveats.}
Two limits of this account deserve emphasis.
First, compression ratio is relative to the compressor: LZ4's window and parsing see local redundancy only, so the band $[0.65, 0.80]$ is a calibrated quantity, not a universal constant, and different compressors, scripts, or domains shift it (Section~\ref{sec:discussion}).
The bilingual example in Figure~\ref{fig:examples} compresses poorly partly for encoding reasons, not only for quality reasons.
Second, the account is first-order: it predicts that a useful interior band exists, not where its endpoints lie.
We therefore determine the endpoints empirically from high-quality reference corpora (Section~\ref{sec:method}).
In short, \textsc{Compel} operationalizes a classical idea: useful training text lives between the poles of redundancy and randomness---compressible enough to be structured, incompressible enough to be informative.
